\begin{table}[htbp]
\centering
\caption{MR and MCL violation rates by SYR2 contaminant reporting status}
\label{tab:syr2_mr_comparison}
\begin{adjustbox}{max width=\textwidth}
\begin{tabular}{lcccc}
\hline\hline
 & \multicolumn{1}{c}{(1)} & \multicolumn{1}{c}{(2)} & \multicolumn{1}{c}{Diff.} & \multicolumn{1}{c}{$p$-value} \\
 & Has SYR2 reading & No SYR2 reading & (1)$-$(2) & \\
\hline
Ever MR violation, 1997--2005 & 32.4\% & 9.8\% & 22.6\%$^{***}$ & 0.000 \\
Ever MR violation, 1985--2005 & 51.1\% & 32.8\% & 18.3\%$^{***}$ & 0.001 \\
Ever MCL violation, 1997--2005 & 1.3\% & 0.0\% & 1.3\%$^{*}$ & 0.083 \\
Ever MCL violation, 1985--2005 & 5.3\% & 1.6\% & 3.7\%$^{*}$ & 0.052 \\
Mean annual violations, 1997--2005 & 0.09 & 0.11 & -0.02$$ & 0.663 \\
\hline
\multicolumn{5}{l}{\textit{N} CWSs: group 1 = 225, group 2 = 122} \\
\multicolumn{5}{l}{\footnotesize Sample: CWSs in the strictly downstream 2SLS sample (downstream HUC, not colocated with a mine).} \\
\multicolumn{5}{l}{\footnotesize \textit{Has SYR2 reading}: $\geq 1$ Six-Year Review contaminant reading for arsenic, nitrate, selenium, chromium, or barium.} \\
\multicolumn{5}{l}{\footnotesize \textit{No SYR2 reading}: CWS is in a state where $\geq 1$ CWS has a SYR2 reading but the CWS itself does not.} \\
\multicolumn{5}{l}{\footnotesize MR violations: any of arsenic, nitrate, or inorganic chemicals monitoring/reporting violation (inorganic chemicals encompasses selenium, chromium, barium).} \\
\multicolumn{5}{l}{\footnotesize MCL violations: any of arsenic, nitrate, or inorganic chemicals maximum contaminant level violation.} \\
\multicolumn{5}{l}{\footnotesize Mean annual violations: sum of MR and MCL violation indicators across arsenic, nitrate, and inorganic chemicals categories (max 6 per CWS-year).} \\
\multicolumn{5}{l}{\footnotesize Diff.~stars from two-sample $t$-test: *** $p{<}0.01$, ** $p{<}0.05$, * $p{<}0.1$.} \\
\hline\hline
\end{tabular}
\end{adjustbox}
\end{table}

